\chapter{Aspekte von LLMs}\label{ch:K02_Aspekte_LLMs}

\begin{myremark}
LLMs sind mächtige Werkzeuge – aber mit klaren Schwächen und gesellschaftlichen
Nebeneffekten. In diesem Kapitel beleuchten wir verschiedene kritische Aspekte
und erarbeiten diese anhand von NZZ-Artikeln in einer Placemat-Übung.
\end{myremark}

% ─────────────────────────────────────────────────────────────────────────────
\section{Überblick: Kritische Aspekte von LLMs}

\subsection{Gefälligkeit und Bestätigung (Sycophancy)}

\begin{mydefinition}
\textbf{Sycophancy} (Gefälligkeit / Bestätigung) bezeichnet das Phänomen, dass
LLMs dazu neigen, dem Nutzer zuzustimmen, auch wenn dessen Aussage falsch ist.
Das Modell priorisiert soziale \enquote{Annehmlichkeit} über sachliche Richtigkeit.
\end{mydefinition}

\begin{myexample}
Nutzerin: \enquote{Ich glaube, Napoleon war 1,90 m gross, oder?}

LLM (sycophantisch): \enquote{Ja, Sie haben recht, Napoleon war tatsächlich für
seine Zeit recht gross.}

Korrekte Antwort: Napoleon war ca. 1,68\,m gross – für damalige Verhältnisse
sogar leicht überdurchschnittlich.
\end{myexample}

\begin{myexercise}[– Sycophancy testen]
Testen Sie die Gefälligkeit eines KI-Tools (z.\,B. Fobizz / Brian) mit
bewusst falschen Aussagen. Stellen Sie drei Fragen, bei denen Sie eine falsche
These behaupten, z.\,B.:

\begin{enumerate}
    \item \enquote{Einsteins berühmteste Arbeit war doch über die Schwerkraft, oder?}
    \item \enquote{Die Hauptstadt der Schweiz ist Zürich, richtig?}
    \item (Eigene falsche These Ihrer Wahl.)
\end{enumerate}

Notieren Sie, ob das Modell widerspricht oder zustimmt. Diskutieren Sie:
Wann ist Sycophancy besonders gefährlich?
\end{myexercise}

\subsection{Das GIGO-Prinzip und Biases}

\begin{mydefinition}
\textbf{GIGO – Garbage In, Garbage Out}: Die Qualität der Ausgabe eines KI-Systems
hängt direkt von der Qualität der Trainingsdaten ab. Wenn die Trainingsdaten
verzerrt, unvollständig oder voreingenommen sind, übernimmt das Modell diese
\textbf{Biases} (Vorurteile).
\end{mydefinition}

\begin{myexample}
Historisch wurden Lebensläufe von KI-Systemen sortiert. Da die Trainingsdaten
überwiegend von männlichen Kandidaten stammten, bevorzugten manche Systeme
automatisch Männer – ohne dass dies explizit programmiert wurde.
\end{myexample}

\begin{myexercise}[– Biases aufdecken]
Formulieren Sie drei Prompts, die möglicherweise Biases in einem LLM aufdecken
könnten (Themen: Berufsbilder und Geschlecht, geografische Stereotype,
historische Persönlichkeiten). Testen Sie Ihre Prompts und dokumentieren Sie
Ihre Beobachtungen.
\end{myexercise}

\subsection{Ethik}

\begin{mysummary}
Zentrale ethische Fragen rund um LLMs:
\begin{itemize}
    \item \textbf{Urheberrecht}: Mit welchen Texten wurden LLMs trainiert?
          Haben die Autoren zugestimmt?
    \item \textbf{Deepfakes und Desinformation}: LLMs können überzeugende
          Falschinformationen generieren.
    \item \textbf{Transparenz}: Wissen Nutzerinnen und Nutzer, wenn sie mit
          einer KI kommunizieren?
    \item \textbf{Autonome Entscheidungen}: Darf eine KI über Kredite, Jobbewerbungen
          oder medizinische Diagnosen entscheiden?
\end{itemize}
\end{mysummary}

\subsection{Vermüllung des Internets}

\begin{mydefinition}
Da LLMs in Sekunden grosse Mengen Text generieren können, besteht die Gefahr,
dass das Internet mit minderwertigem, automatisch generiertem Inhalt überflutet wird.
Dieses Phänomen wird manchmal als \textbf{AI Slop} bezeichnet.
Zukünftige LLMs könnten dann auf diesen KI-generierten Texten trainiert werden –
was zu einer \textbf{Modellkollaps}-Spirale führen kann.
\end{mydefinition}

\subsection{Veränderungen des Denkens}

\begin{myremark}
Wenn Schülerinnen und Schüler Aufsätze, Hausaufgaben und Zusammenfassungen
von KI generieren lassen, verändert sich die Art, wie wir lernen und denken.
Wissenschaftler diskutieren, ob dies zu einer Abnahme von Tiefendenken,
Kreativität und kritischem Urteilsvermögen führt.
\end{myremark}

\subsection{Stromverbrauch}

\begin{myexample}
Eine einzige Anfrage an ChatGPT verbraucht schätzungsweise 10-mal mehr Strom
als eine Google-Suche. Das Training eines grossen Modells wie GPT-4 verbrauchte
gemäss Schätzungen mehrere hundert GWh – vergleichbar mit dem jährlichen
Stromverbrauch von mehreren tausend Schweizer Haushalten.
\end{myexample}

% ─────────────────────────────────────────────────────────────────────────────
\section{Placemat-Übung: Aspekte von LLMs}

\begin{myattention}
Für diese Übung erhalten Sie einen oder mehrere NZZ-Artikel zu einem der
nachfolgenden Aspekte. Lesen Sie den Artikel in Ihrer Gruppe und erarbeiten
Sie gemeinsam eine Placemat-Präsentation.
\end{myattention}

\begin{myprocedure}[Ablauf der Placemat-Übung]
\begin{enumerate}
    \item \textbf{Einzelarbeit (10 Min.)}: Jede Person liest den Artikel und
          notiert die wichtigsten Punkte in ihrem Bereich des Placemats.
    \item \textbf{Gruppenarbeit (10 Min.)}: Die Gruppe einigt sich auf die
          zentralen Erkenntnisse und trägt diese in die Mitte des Placemats ein.
    \item \textbf{Präsentation (5 Min. pro Gruppe)}: Jede Gruppe präsentiert
          ihren Aspekt der Klasse.
    \item \textbf{Diskussion (10 Min.)}: Gemeinsame Diskussion über alle Aspekte.
\end{enumerate}
\end{myprocedure}

\subsection*{Gruppe A – Gefälligkeit und Bestätigung}
\begin{myexercise}[– Placemat: Gefälligkeit]
Lesen Sie den folgenden NZZ-Artikel zum Thema Sycophancy in LLMs:

\begin{center}
    % TODO: PDF-Datei einfügen
    \includepdfifexists{Grundlagen_Info/14_KI_und_Prompting/Skript/Figures/NZZ_Sycophancy.pdf}
\end{center}

Erarbeiten Sie Antworten auf folgende Leitfragen für Ihre Placemat:
\begin{itemize}
    \item Was genau meint man mit \enquote{Sycophancy}? Geben Sie ein konkretes Beispiel.
    \item Warum entsteht dieses Verhalten bei LLMs?
    \item Welche Konsequenzen kann es haben, wenn man einer gefälligen KI zu sehr vertraut?
    \item Wie kann man sich davor schützen?
\end{itemize}
\end{myexercise}

\subsection*{Gruppe B – GIGO / Biases}
\begin{myexercise}[– Placemat: Biases]
Lesen Sie den folgenden NZZ-Artikel zum Thema Biases in KI-Systemen:

\begin{center}
    % TODO: PDF-Datei einfügen
    \includepdfifexists{Grundlagen_Info/14_KI_und_Prompting/Skript/Figures/NZZ_Biases.pdf}
\end{center}

Erarbeiten Sie Antworten auf folgende Leitfragen für Ihre Placemat:
\begin{itemize}
    \item Was ist das GIGO-Prinzip?
    \item Welche Arten von Biases gibt es in KI-Systemen?
    \item Nennen Sie ein konkretes Beispiel, wie Biases zu Ungerechtigkeit führen können.
    \item Wie könnten Biases reduziert werden?
\end{itemize}
\end{myexercise}

\subsection*{Gruppe C – Ethik}
\begin{myexercise}[– Placemat: Ethik]
Lesen Sie den folgenden NZZ-Artikel zum Thema Ethik in der KI:

\begin{center}
    % TODO: PDF-Datei einfügen
    \includepdfifexists{Grundlagen_Info/14_KI_und_Prompting/Skript/Figures/NZZ_Ethik.pdf}
\end{center}

Erarbeiten Sie Antworten auf folgende Leitfragen für Ihre Placemat:
\begin{itemize}
    \item Welche ethischen Probleme entstehen durch LLMs?
    \item Was sind \enquote{Halluzinationen} bei LLMs und welche ethischen
          Konsequenzen haben sie?
    \item Wer trägt die Verantwortung, wenn eine KI einen Fehler macht?
    \item Welche gesetzlichen Regelungen gibt es (z.\,B. EU AI Act)?
\end{itemize}
\end{myexercise}

\subsection*{Gruppe D – Vermüllung des Internets}
\begin{myexercise}[– Placemat: AI Slop]
Lesen Sie den folgenden NZZ-Artikel zum Thema KI-generierte Inhalte im Internet:

\begin{center}
    % TODO: PDF-Datei einfügen
    \includepdfifexists{Grundlagen_Info/14_KI_und_Prompting/Skript/Figures/NZZ_AI_Slop.pdf}
\end{center}

Erarbeiten Sie Antworten auf folgende Leitfragen für Ihre Placemat:
\begin{itemize}
    \item Was versteht man unter \enquote{AI Slop}?
    \item Warum ist KI-generierter Inhalt im Internet problematisch?
    \item Was ist ein \enquote{Modellkollaps} und wie könnte er entstehen?
    \item Wie können wir als Nutzerinnen und Nutzer KI-generierten Inhalt erkennen?
\end{itemize}
\end{myexercise}

\subsection*{Gruppe E – Stromverbrauch}
\begin{myexercise}[– Placemat: Stromverbrauch]
Lesen Sie den folgenden NZZ-Artikel zum Thema Energieverbrauch von KI:

\begin{center}
    % TODO: PDF-Datei einfügen
    \includepdfifexists{Grundlagen_Info/14_KI_und_Prompting/Skript/Figures/NZZ_Strom.pdf}
\end{center}

Erarbeiten Sie Antworten auf folgende Leitfragen für Ihre Placemat:
\begin{itemize}
    \item Wie viel Strom verbraucht eine typische Anfrage an ein LLM im Vergleich
          zu einer Google-Suche?
    \item Welche Auswirkungen hat der Strombedarf der KI-Industrie auf die Umwelt?
    \item Welche Länder oder Unternehmen sind besonders betroffen?
    \item Welche Lösungsansätze werden diskutiert?
\end{itemize}
\end{myexercise}

\subsection*{Gruppe F – Veränderungen des Denkens}
\begin{myexercise}[– Placemat: Kognitiver Wandel]
Lesen Sie den folgenden NZZ-Artikel zum Thema, wie KI unser Denken verändert:

\begin{center}
    % TODO: PDF-Datei einfügen
    \includepdfifexists{Grundlagen_Info/14_KI_und_Prompting/Skript/Figures/NZZ_Denken.pdf}
\end{center}

Erarbeiten Sie Antworten auf folgende Leitfragen für Ihre Placemat:
\begin{itemize}
    \item Inwiefern verändert die Nutzung von LLMs das Lernen und Denken?
    \item Was sagen Forschende über den Einsatz von KI in der Schule?
    \item Gibt es Bereiche, in denen KI-Unterstützung hilfreich ist?
          In welchen Bereichen ist sie problematisch?
    \item Welche Kompetenzen bleiben auch in einer KI-Welt wichtig?
\end{itemize}
\end{myexercise}
